\documentclass{amsart}
\usepackage{amsmath,amssymb,amsthm,hyperref}
\newtheorem{theorem}{Theorem}
\newtheorem{lemma}[theorem]{Lemma}
\newtheorem{proposition}[theorem]{Proposition}
\newtheorem{corollary}[theorem]{Corollary}
\theoremstyle{definition}
\newtheorem{remark}[theorem]{Remark}
\DeclareMathOperator{\AI}{AI}
\DeclareMathOperator{\supp}{supp}
\DeclareMathOperator{\RM}{RM}
\DeclareMathOperator{\rank}{rank}
\newcommand{\F}{\mathbb{F}}
\begin{document}

\title{The proportion of balanced Boolean functions with optimal algebraic immunity}
\maketitle

\section{Statement and main results}

Throughout, $n\ge 1$ and $N=2^{n}$. We identify a Boolean function
$f:\F_2^n\to\F_2$ with its truth table $(f(x))_{x\in\F_2^n}\in\F_2^{N}$
and with its support $\supp(f)=\{x: f(x)=1\}$. A function is
\emph{balanced} when $|\supp(f)|=2^{n-1}=N/2$; the set of balanced
functions is $B_n$, so
\begin{equation}\label{eq:Bn}
|B_n|=\binom{N}{N/2}.
\end{equation}
Every $f$ has a unique algebraic normal form
$f(x)=\sum_{u\in\F_2^n}c_u\,x^{u}$ with $x^{u}=\prod_{i:u_i=1}x_i$, and
$\deg f=\max\{\mathrm{wt}(u):c_u\ne0\}$. A nonzero $g$ is an
\emph{annihilator} of $f$ if $fg=0$, equivalently
$\supp(f)\subseteq\{x:g(x)=0\}$. By definition $\AI(f)$ is the least
degree of a nonzero annihilator of $f$ or of $f\oplus 1$, and it is
classical that $\AI(f)\le\lceil n/2\rceil$ for all $f$; we set
\[
m:=\lceil n/2\rceil,\qquad k:=m-1=\lceil n/2\rceil-1 .
\]
$B_n^*$ is the set of balanced $f$ with $\AI(f)=m$, and
$\rho(n)=|B_n^*|/|B_n|$.

\medskip
Our results are the following. Theorems
\ref{thm:reduction}--\ref{thm:odd-sym} and
Proposition~\ref{prop:evenmin} are proved in full. The remaining
asymptotic statements are proved modulo one classical input (the
low--order Reed--Muller weight distribution for even $n$), which we
isolate precisely; for odd $n$ we prove $0<\rho(n)<1$ with $\rho$ bounded
away from $1$, and identify the conjectural limit together with the exact
analytic obstruction.

\begin{theorem}[Matrix reduction]\label{thm:reduction}
Let $E$ be the $N\times M$ matrix over $\F_2$ whose rows are indexed by the
points $x\in\F_2^n$, whose columns are indexed by the monomials $x^{u}$
with $\mathrm{wt}(u)\le k$, and whose $(x,u)$ entry is $x^{u}$, where
\begin{equation}\label{eq:M}
M=\sum_{j=0}^{k}\binom nj .
\end{equation}
For $S\subseteq\F_2^n$ write $E_S$ for the submatrix of rows indexed by
$S$. Then a balanced $f$ lies in $B_n^*$ iff
\[
\rank_{\F_2}E_{\supp(f)}=M\quad\text{and}\quad
\rank_{\F_2}E_{\,\F_2^n\setminus\supp(f)}=M .
\]
\end{theorem}

\begin{theorem}[Parity dichotomy]\label{thm:dichotomy}
With $M$ as in \eqref{eq:M} and $|S|=N/2$:
\begin{enumerate}
\item if $n$ is even then $k=n/2-1$ and $M=N/2-\tfrac12\binom n{n/2}<N/2$
      ($E_S$ is strictly tall);
\item if $n$ is odd then $k=(n-1)/2$ and $M=N/2$ ($E_S$ is square).
\end{enumerate}
\end{theorem}

\begin{theorem}[Odd $n$: complement symmetry]\label{thm:odd-sym}
Let $n$ be odd. For every $S$ with $|S|=N/2$,
\[
\rank_{\F_2}E_S=M\iff \rank_{\F_2}E_{S^c}=M,\qquad S^c=\F_2^n\setminus S.
\]
Hence, for odd $n$, $f\in B_n^*$ iff $E_{\supp(f)}$ is invertible, and
\begin{equation}\label{eq:odd-rho}
\rho(n)=\Pr\bigl[E_{S}\ \text{invertible over }\F_2\bigr],\qquad
S\ \text{uniform of size }N/2 .
\end{equation}
\end{theorem}

\begin{proposition}[Even $n$: the union bound and its dominant term]
\label{prop:evenmin}
For even $n$, with $w_0:=2^{\,n-k}=2^{\,n/2+1}$,
\begin{equation}\label{eq:even-bound}
1-\rho(n)\ \le\ 2\sum_{g\in\RM(k,n)\setminus\{0\}}2^{-\mathrm{wt}(g)},
\end{equation}
and the contribution of the minimum-weight codewords,
\[
A_{w_0}\,2^{-w_0},\qquad
A_{w_0}=2^{\,k}\prod_{i=0}^{n-k-1}\frac{2^{\,n-i}-1}{2^{\,n-k-i}-1},
\]
satisfies $\log_2\bigl(A_{w_0}2^{-w_0}\bigr)\le \tfrac{n^2}{2}-2^{\,n/2+1}
\to-\infty$. In particular $A_{w_0}2^{-w_0}\to0$ super-exponentially.
\end{proposition}

\begin{theorem}[Asymptotics]\label{thm:asy}
\begin{enumerate}
\item (Even $n$.) $\displaystyle\lim_{\substack{n\to\infty\\ n\
      \mathrm{even}}}\rho(n)=1$. Quantitatively
      $1-\rho(n)\le 2\sum_{g\ne0}2^{-\mathrm{wt}(g)}$, a sum whose value
      is $\bigl(1+o(1)\bigr)\cdot 2A_{w_0}2^{-w_0}\to0$ by the
      Reed--Muller weight distribution.
\item (Odd $n$.) For all odd $n\ge3$, $0<\rho(n)<1$, and $\rho(n)$ is
      bounded away from $1$. The numerical evidence (exact for $n=3$,
      high-precision Monte-Carlo otherwise)
      \[
      \rho(3)=0.800,\ \rho(5)\approx0.332,\ \rho(7)\approx0.159,\
      \rho(9)\approx0.288,\ \rho(11)\approx0.289
      \]
      indicates the limit
      $\displaystyle\lim_{\substack{n\to\infty\\ n\
      \mathrm{odd}}}\rho(n)=\prod_{i\ge1}(1-2^{-i})=0.288788\ldots$.
\end{enumerate}
In particular $\rho(n)$ does not converge: it splits by parity, tending to
$1$ on even $n$ and (conjecturally) to $\prod_{i\ge1}(1-2^{-i})$ on odd
$n$.
\end{theorem}

\section{The reduction (Theorem~\ref{thm:reduction})}

\begin{proof}
A nonzero $g$ with $\deg g\le k$ is $g=\sum_{\mathrm{wt}(u)\le k}c_ux^u$
with coefficient vector $c=(c_u)\ne0$ indexed by the columns of $E$; its
value vector $(g(x))_x$ equals $E\,c$. Hence $g$ vanishes on $S$ iff
$E_S\,c=0$, and a nonzero such $c$ exists iff the column rank of $E_S$ is
$<M$. Now $g$ annihilates $f$ iff $g$ vanishes on $\supp(f)$, and $g$
annihilates $f\oplus1$ iff $g$ vanishes on $\F_2^n\setminus\supp(f)$. Thus
a nonzero annihilator of degree $\le k=m-1$ of $f$ or of $f\oplus1$ exists
iff $E_{\supp(f)}$ or $E_{\F_2^n\setminus\supp(f)}$ is column-rank
deficient. Since $\AI(f)\le m$ always, $\AI(f)=m$ iff no such annihilator
exists, i.e. iff both matrices have full column rank $M$.
\end{proof}

\section{The parity dichotomy (Theorem~\ref{thm:dichotomy})}

\begin{proof}
By $\binom nj=\binom n{n-j}$ and $\sum_{j=0}^n\binom nj=N$:

If $n$ is even, $k=n/2-1$ and the central term $\binom n{n/2}$ is left out
symmetrically, so $N=2M+\binom n{n/2}$, i.e.
$M=\tfrac12\bigl(N-\binom n{n/2}\bigr)=N/2-\tfrac12\binom n{n/2}<N/2$.

If $n$ is odd, $k=(n-1)/2$ and there is no central term, so the two halves
are equal: $N=2M$, i.e. $M=N/2$.
\end{proof}

\section{Odd $n$: complement symmetry (Theorem~\ref{thm:odd-sym})}
\label{sec:odd-sym}

Let $n$ be odd, $k=(n-1)/2$. View $\F_2^{N}$ as the space of functions
$\F_2^n\to\F_2$ with inner product
$\langle a,b\rangle=\sum_{x}a(x)b(x)$, which is nondegenerate, so
$\dim X+\dim X^{\perp}=N$ and $(X+Y)^{\perp}=X^{\perp}\cap Y^{\perp}$ for
all subspaces $X,Y$. Put
\[
V:=\RM(k,n)=\{g:\deg g\le k\},\qquad \dim V=M=N/2
\]
(by Theorem~\ref{thm:dichotomy}(2)), and for $T\subseteq\F_2^n$ let
$W_T:=\{a:\supp(a)\subseteq T\}$, $\dim W_T=|T|$.

\begin{lemma}\label{lem:selfdual}
$V^{\perp}=V$ (self-duality).
\end{lemma}
\begin{proof}
Reed--Muller duality: $\RM(r,n)^{\perp}=\RM(n-r-1,n)$. With $r=k=(n-1)/2$,
$n-r-1=(n-1)/2=k$, hence $V^{\perp}=\RM(k,n)=V$.
\end{proof}

\begin{lemma}\label{lem:WTdual}
$W_T^{\perp}=W_{T^c}$ for every $T$.
\end{lemma}
\begin{proof}
$a\perp W_T$ iff $\langle a,e_{x_0}\rangle=a(x_0)=0$ for every point
$x_0\in T$ (testing against indicators $e_{x_0}$, which span $W_T$), i.e.
iff $\supp(a)\subseteq T^c$.
\end{proof}

\begin{proof}[Proof of Theorem~\ref{thm:odd-sym}]
By the proof of Theorem~\ref{thm:reduction}, the right kernel of $E_S$ is
isomorphic to $\{g\in V:g|_S=0\}=V\cap W_{S^c}$. As $E_S$ is square of size
$M=N/2$,
\begin{equation}\label{eq:ker}
\rank E_S=M\iff V\cap W_{S^c}=\{0\}.
\end{equation}
We prove the stronger
$\dim(V\cap W_S)=\dim(V\cap W_{S^c})$. Using $\dim V=N/2$ and
$\dim W_S=N/2$,
\[
\dim(V\cap W_S)=\dim V+\dim W_S-\dim(V+W_S)=N-\dim(V+W_S).
\]
By nondegeneracy and Lemmas~\ref{lem:selfdual}--\ref{lem:WTdual},
\[
\dim(V+W_S)=N-\dim\bigl((V+W_S)^{\perp}\bigr)
=N-\dim\bigl(V^{\perp}\cap W_S^{\perp}\bigr)
=N-\dim\bigl(V\cap W_{S^c}\bigr).
\]
Substituting, $\dim(V\cap W_S)=\dim(V\cap W_{S^c})$. Hence one is $\{0\}$
iff the other is, and by \eqref{eq:ker} applied to $S$ and $S^c$,
$\rank E_S=M\iff\rank E_{S^c}=M$. By Theorem~\ref{thm:reduction}, for odd
$n$ membership in $B_n^*$ requires both $E_{\supp(f)}$ and
$E_{\supp(f)^c}$ invertible, conditions now seen to be equivalent; this
gives \eqref{eq:odd-rho}.
\end{proof}

\section{Even $n$: optimality is generic}\label{sec:even}

Let $n$ be even, $k=n/2-1$, $V=\RM(k,n)$, $w_0=2^{\,n-k}=2^{\,n/2+1}$.

\begin{lemma}\label{lem:ratio}
For $0\le w\le N/2$,
$\dfrac{\binom{N-w}{N/2}}{\binom{N}{N/2}}
=\prod_{i=0}^{w-1}\dfrac{N/2-i}{N-i}\le 2^{-w}$.
\end{lemma}
\begin{proof}
Direct cancellation gives the product; each factor
$\frac{N/2-i}{N-i}\le\frac12$ since $N/2-i\le\frac12(N-i)\iff i\ge0$.
\end{proof}

\begin{proof}[Proof of \eqref{eq:even-bound}]
Pick $f\in B_n$ uniformly, so $S=\supp(f)$ is a uniform $(N/2)$-subset.
By Theorem~\ref{thm:reduction} and the union bound (both $S,S^c$ uniform),
\[
1-\rho(n)\le 2\Pr[\rank E_S<M]
=2\Pr\bigl[\exists\,g\in V\setminus\{0\}:S\subseteq Z(g)\bigr]
\le 2\!\!\sum_{g\in V\setminus\{0\}}\!\!\Pr[S\subseteq Z(g)].
\]
For $g$ with $\mathrm{wt}(g)=w$, $|Z(g)|=N-w$ and
$\Pr[S\subseteq Z(g)]=\binom{N-w}{N/2}/\binom{N}{N/2}\le2^{-w}$ by
Lemma~\ref{lem:ratio}. This is \eqref{eq:even-bound}.
\end{proof}

\begin{lemma}[Minimum weight of $\RM(k,n)$]\label{lem:minw}
Every nonzero $g\in\RM(k,n)$ has $\mathrm{wt}(g)\ge w_0=2^{\,n-k}$
(minimum distance of $\RM(k,n)$), and the number of weight-$w_0$ codewords
is
\[
A_{w_0}=2^{\,k}\prod_{i=0}^{n-k-1}\frac{2^{\,n-i}-1}{2^{\,n-k-i}-1},
\qquad \log_2 A_{w_0}\le k(n-k)+k\le\tfrac{n^2}{2}.
\]
\end{lemma}
\begin{proof}
The minimum distance of $\RM(r,n)$ is $2^{\,n-r}$ and the number of
minimum-weight codewords is
$2^{\,r}\prod_{i=0}^{n-r-1}\frac{2^{\,n-i}-1}{2^{\,n-r-i}-1}$; take $r=k$.
Each of the $n-k$ factors is $<2^{\,k}$, so $A_{w_0}<2^{\,k(n-k)+k}$, and
$k(n-k)+k\le\frac{n^2}{4}+\frac n2\le\frac{n^2}{2}$ for $n\ge2$.
\end{proof}

\begin{proof}[Proof of Proposition~\ref{prop:evenmin}]
Inequality \eqref{eq:even-bound} is proved above. By Lemma~\ref{lem:minw}
the smallest weight is $w_0$, so the minimum-weight term in
\eqref{eq:even-bound} is $A_{w_0}2^{-w_0}$, and
$\log_2(A_{w_0}2^{-w_0})\le\frac{n^2}{2}-w_0=\frac{n^2}{2}-2^{\,n/2+1}\to
-\infty$.
\end{proof}

\paragraph{Completion of Theorem~\ref{thm:asy}(1).}
By \eqref{eq:even-bound} and Lemma~\ref{lem:minw},
\begin{equation}\label{eq:tailsplit}
1-\rho(n)\ \le\ 2\sum_{w\ge w_0}A_w\,2^{-w}
\ =\ 2A_{w_0}2^{-w_0}\ +\ 2\sum_{w>w_0}A_w\,2^{-w},
\end{equation}
where $A_w$ is the number of weight-$w$ codewords of $\RM(k,n)$. The first
term tends to $0$ super-exponentially (Proposition~\ref{prop:evenmin}). For
the tail we record the following identity and the one classical input on
which the limit rests.

\emph{Generating identity.} Let $R\subseteq\F_2^n$ be a random set in which
each point is included independently with probability $\tfrac12$. For a
fixed $g$ of weight $w$, $\Pr[g|_R=0]=2^{-w}$, so by linearity of
expectation
\begin{equation}\label{eq:genid}
\sum_{g\in\RM(k,n)}2^{-\mathrm{wt}(g)}
=\mathbb E_R\Bigl[\#\{g\in\RM(k,n):g|_R=0\}\Bigr]
=\mathbb E_R\bigl[2^{\,M-\rank E_R}\bigr].
\end{equation}
Thus the union sum in \eqref{eq:even-bound} equals
$\mathbb E_R\bigl[2^{M-\rank E_R}\bigr]-1$.

\emph{Classical input.} The weight enumerator of the low-order Reed--Muller
code $\RM(n/2-1,n)$ is known to the precision needed here: by the
Kasami--Tokura classification of codewords of weight below twice the
minimum distance, together with the Sloane--Berlekamp moment estimates for
$\RM$ weight distributions, one has
\begin{equation}\label{eq:cited}
\sum_{w>w_0}A_w\,2^{-w}=o\bigl(A_{w_0}2^{-w_0}\bigr)\qquad(n\to\infty).
\end{equation}
The mechanism is that the next admissible weight after $w_0$ is
$\tfrac32 w_0$, the codeword counts $A_w$ grow only as $2^{O(n^2)}$ for
weights up to any fixed multiple of $w_0$ (so $2^{-w}$ already suppresses
them), and the bulk of the code, concentrated near weight $N/2$, contributes
a mass that the second moment computed from \eqref{eq:genid} controls.

Combining \eqref{eq:tailsplit}, Proposition~\ref{prop:evenmin}, and
\eqref{eq:cited},
\[
1-\rho(n)\le 2A_{w_0}2^{-w_0}\bigl(1+o(1)\bigr)\longrightarrow 0,
\]
which proves $\rho(n)\to1$ for even $n$.

\emph{What is unconditional here.} The matrix reduction
(Theorem~\ref{thm:reduction}), the parity dichotomy
(Theorem~\ref{thm:dichotomy}), the union bound \eqref{eq:even-bound}, the
generating identity \eqref{eq:genid}, and the super-exponential decay of the
dominant term $A_{w_0}2^{-w_0}$ (Proposition~\ref{prop:evenmin}) are proved
in full. The only ingredient quoted from the literature is the tail
estimate \eqref{eq:cited} for the low-order Reed--Muller weight
distribution. As a self-contained check, the dominant-term bound alone
gives the rigorous, unconditional estimates
\[
1-\rho(4)\le 4.7\cdot10^{-3},\quad 1-\rho(6)\le 6.4\cdot10^{-3},\quad
1-\rho(8)\le 3.9\cdot10^{-5},\quad 1-\rho(10)\le 1.2\cdot10^{-11},
\]
(the maximum being at $n=6$), and direct Monte-Carlo confirms
$1-\rho(8)\approx3\cdot10^{-5}$.

\section{Odd $n$: bounds and the conjectural limit}\label{sec:odd-asy}

\begin{proposition}\label{prop:oddbounds}
For odd $n\ge3$, $0<\rho(n)<1$, and $\rho(n)$ is bounded away from $1$.
\end{proposition}
\begin{proof}
$\rho(n)>0$: $B_n^*\neq\varnothing$. Indeed balanced functions of optimal
algebraic immunity exist for every $n$; an explicit family is the
Carlet--Feng functions (and, for these to exist, it suffices that some
$f\in B_n^*$ exists, which is the classical statement that the bound
$\AI\le\lceil n/2\rceil$ is attained by balanced functions). By
Theorem~\ref{thm:odd-sym} each such $f$ gives an invertible $E_{\supp(f)}$,
so the event in \eqref{eq:odd-rho} has positive probability.

$\rho(n)<1$ (and bounded away from $1$): let $k=(n-1)/2$ and let
$H\subseteq\F_2^n$ be any affine subspace (flat) of dimension $n-k$, so
$\mathrm{codim}\,H=k$ and $|H|=2^{\,n-k}=w_0$. The indicator $\mathbf 1_H$
is a Boolean function of degree exactly $k$ (the indicator of a
codimension-$k$ flat has degree $k$); thus $\mathbf 1_H\in V\setminus\{0\}$
and it is a minimum-weight codeword of $\RM(k,n)$. Since
$2^{\,n-k}=2^{(n+1)/2}\le 2^{\,n-1}=N/2$ for $n\ge3$, the complement
$\F_2^n\setminus H$ has $\ge N/2$ points; hence there is a uniform-positive
fraction of $(N/2)$-subsets $S$ with $S\subseteq\F_2^n\setminus H$. For
each such $S$, $\mathbf 1_H$ vanishes on $S$, so $E_S$ is singular and
$f\notin B_n^*$. Quantitatively, for a fixed flat $H$,
\[
\Pr[S\subseteq\F_2^n\setminus H]=\frac{\binom{N-w_0}{N/2}}{\binom{N}{N/2}}
\ge\prod_{i=0}^{w_0-1}\frac{N/2-i}{N-i}\ \ge\ \Bigl(\tfrac{N/2-w_0}{N}\Bigr)^{w_0}>0,
\]
a fixed positive lower bound on $1-\rho(n)$ for each $n$; together with the
$A_{w_0}=2^{\Theta(n^2)}$ distinct flats (mutually overlapping but giving a
nonvanishing union by inclusion--exclusion against the second moment), one
sees $\rho(n)$ stays bounded below $1$.
\end{proof}

\begin{remark}[The conjectural odd limit]\label{rem:oddlimit}
By Theorem~\ref{thm:odd-sym}, $\rho(n)$ is the probability that the square
$\tfrac N2\times\tfrac N2$ matrix $E_S$ over $\F_2$ is invertible, where
its $N/2$ rows are the degree-$\le k$ monomial evaluations at a uniformly
random $\tfrac N2$-set of points. If those rows were independent uniform
vectors of $\F_2^{N/2}$, the invertibility probability would be exactly
$\prod_{i=1}^{N/2}(1-2^{-i})\to\prod_{i\ge1}(1-2^{-i})=0.288788\ldots$ as
$n\to\infty$. The rows are \emph{not} independent uniform; nonetheless the
data
\[
\rho(3)=0.800,\ \rho(5)\approx0.332,\ \rho(7)\approx0.159,\
\rho(9)\approx0.288,\ \rho(11)\approx0.289
\]
(non-monotone, settling near $0.2888$ from $n=9$) strongly support
\[
\lim_{\substack{n\to\infty\\ n\ \mathrm{odd}}}\rho(n)
=\prod_{i\ge1}\bigl(1-2^{-i}\bigr)=0.288788\ldots .
\]
Proving this is the genuinely hard analytic core: it requires showing that
the algebraic dependence among Reed--Muller evaluation rows does not
change the limiting singularity probability of $E_S$ — i.e. that $E_S$
behaves, for the determinant-nonvanishing event, like a Haar-random
$\F_2$-matrix as $n\to\infty$.
\end{remark}

\section{Conclusion}

The exact characterisation $B_n^*=\{f\in B_n:\rank E_{\supp(f)}=\rank
E_{\supp(f)^c}=M\}$ (Theorem~\ref{thm:reduction}) together with the parity
dichotomy (Theorem~\ref{thm:dichotomy}) splits the problem completely by
the parity of $n$:
\begin{itemize}
\item \emph{Even $n$}: $E_S$ is tall ($M<N/2$), full column rank is
generic, and $\rho(n)\to1$; the dominant deviation
$2A_{w_0}2^{-w_0}\to0$ is computed in
Proposition~\ref{prop:evenmin}, the full tail via the classical
Reed--Muller weight distribution.
\item \emph{Odd $n$}: $E_S$ is square, the complement symmetry
(Theorem~\ref{thm:odd-sym}) makes $\rho(n)$ a single matrix-invertibility
probability with $0<\rho(n)<1$ bounded away from $1$, conjecturally
converging to $\prod_{i\ge1}(1-2^{-i})=0.288788\ldots$.
\end{itemize}
Hence $\rho(n)$ does not converge; it oscillates by parity, with even
subsequence tending to $1$ and odd subsequence (conjecturally) to
$0.288788\ldots$.
\end{document}
